\section{External reference qualification}
The Greater Wellington intake returned records for nine selected sites. Under the originally specified quality-code-600 criterion no complete reference hot year was qualified; the provider describes code 200 as unverified, while other returned legacy code meanings remain unresolved. A NEON PRT archive supplied half-hourly records at 24 river or stream sites. With one fixed location code per site, no site supplied four complete years under either 48 quality-passing intervals per UTC day or the 46-of-48 sensitivity. The latter is an observed-interval average rather than a complete physical day; a location code does not establish unchanged hardware.

After the original intake and quality screen, we separately tested an optimistic time-support bound on the same nine cached Greater Wellington records: all quality flags were ignored, while the original fixed-NZST cadence and fully observed day rule were retained and finite values required. Across these records there were 10 potentially complete July--June years, with zero to two per site; no site reached four. This post-hoc feasibility screen is conditional on the cached series and these completeness rules, not a quality assessment, reference qualification, or recovery score. The annual counts were recomputed independently from the same cached XML using a separate day-count implementation; the XML parser and source observations were shared.

An initial 15-station USGS intake had already been processed in the national panel, and ten stations had been scored in an earlier external comparison. The subsequent eight-station roster excluded these and other identified source, evaluation and cached-value contacts; no station passed the four-year reference requirement after the documented quality checks. A further five metadata candidates with prior network screening had no four years of complete approved daily counts in 2010--2024; counts can include estimated values and do not establish subdaily quality. Five selected R\'eunion Hub'Eau sites provided records with quality code 1, but no site supplied four complete July--June years under the 24-distinct-valid-hour rule; the provider-day time zone is unresolved. The same intake results do not describe the fraction of eligible stations in any complete network.

At ECCC Red River at Emerson, the USGS alias and an upstream source station place the record on a previously encountered river. The files use the word VALIDATED and, for some rows, a station identifier different from the catalog; the mapping to literal V and numeric quality grades and the daylight-saving convention are unresolved. Under the fixed literal V and grade 1/2 rule, zero complete local-clock years qualified. These exclusions are not an assertion that the rejected observations are erroneous; no numerical recovery score was calculated for this site. Data-source citations appear in the main article.
